\documentclass[12pt, a4paper]{article}

\usepackage[utf8]{inputenc}
\usepackage[T1]{fontenc}
\usepackage[francais]{babel}
\usepackage{graphicx}
\usepackage{fontawesome}

\usepackage{helvet}
\renewcommand{\familydefault}{\sfdefault}


\begin{document}


%%%%%%%%%%%%%%%%%%%%%%%%%%%%%%%%%%%%%%%%%%%%%%%%%%%%%
%Couverture:
%%%%%%%%%%%%%%%%%%%%%%%%%%%%%%%%%%%%%%%%%%%%%%%%%%%%%
\begin{titlepage}
  \begin{center}

    % Title
    \rule{\linewidth}{0.5mm} \\[0.4cm]
    { \huge \bfseries {\LARGE{Rapport de projet}}
    \\Retouche d'images sur Android\\[0.4cm] }
    \rule{\linewidth}{0.5mm} \\[1.5cm]
    
    {\LARGE {\textbf{TODO noms}}}\\[0.5cm]
    {\LARGE L3 Informatique}\\[1.5cm]
    
    {\Large Projet technologique}\\[0.5cm]
    {\Large 2020}\\[1.5cm]
    
    \includegraphics[width=1\textwidth]{report_src/logoUB.jpg}
    
    
  \end{center}
\end{titlepage}



\tableofcontents
\clearpage 


%%%%%%%%%%%%%%%%%%%%%%%%%%%%%%%%%%%%%%%%%%%%%%%%%%%%%
% Introduction :
%%%%%%%%%%%%%%%%%%%%%%%%%%%%%%%%%%%%%%%%%%%%%%%%%%%%%
\section{Introduction :} \label{intro}
TODO

 
\clearpage

%%%%%%%%%%%%%%%%%%%%%%%%%%%%%%%%%%%%%%%%%%%%%%%%%%%%%
%L'application :
%%%%%%%%%%%%%%%%%%%%%%%%%%%%%%%%%%%%%%%%%%%%%%%%%%%%%
\section{L'application :}

\subsection{Lancement :}
TODO


\subsection{Utilisation :}
TODO



%%%%%%%%%%%%%%%%%%%%%%%%%%%%%%%%%%%%%%%%%%%%%%%%%%%%%
% Effets disponibles :
%%%%%%%%%%%%%%%%%%%%%%%%%%%%%%%%%%%%%%%%%%%%%%%%%%%%%
\section{Effets :}
TODO




%%%%%%%%%%%%%%%%%%%%%%%%%%%%%%%%%%%%%%%%%%%%%%%%%%%%%
%Structure du projet:
%%%%%%%%%%%%%%%%%%%%%%%%%%%%%%%%%%%%%%%%%%%%%%%%%%%%%
\section{Structure du projet :}
\subsection{Structure graphique Android et navigation :}
Pour afficher certains éléments d'interface et notamment les informations de l'image (bouton \faInfoCircle) nous utilisons des \textbf{Fragment}. En effet une activité supplémentaire n'est pas nécessaire car cette petite partie de l'application ne correspond pas à un point d'entrée de l'application. Par ailleurs changer de fragment (plutôt que de changer directement de layout) pourrait faciliter l'implémentation d'une interface différente, pour tablette par exemple.
\\
On notera que la structure actuelle de l'application n'exploite pas encore comme il faut les fragments, l'essentiel du code de l'interface est toujours directement utilisé depuis l'activité principale.
\\


%%%%%%%%%%%%%%%%%%%%%%%%%%%%%%%%%%%%%%%%%%%%%%%%%%%%%
% Tests de performance :
%%%%%%%%%%%%%%%%%%%%%%%%%%%%%%%%%%%%%%%%%%%%%%%%%%%%%
\section{Performances :}

\subsection{Temps d'exécution :}
TODO






\subsection{Mémoire :}
TODO


%%%%%%%%%%%%%%%%%%%%%%%%%%%%%%%%%%%%%%%%%%%%%%%%%%%%%
% Remarques et améliorations :
%%%%%%%%%%%%%%%%%%%%%%%%%%%%%%%%%%%%%%%%%%%%%%%%%%%%%
\section{Remarques et améliorations  :}

\subsection{Remarques sur le code}
TODO

\subsection{Améliorations à court terme :}
TODO


\clearpage

%%%%%%%%%%%%%%%%%%%%%%%%%%%%%%%%%%%%%%%%%%%%%%%%%%%%%
% Conclusion :
%%%%%%%%%%%%%%%%%%%%%%%%%%%%%%%%%%%%%%%%%%%%%%%%%%%%%
\section{Conclusion :}
TODO



%%%%%%%%%%%%%%%%%%%%%%%%%%%%%%%%%%%%%%%%%%%%%%%%%%%%%
% Annexes :
%%%%%%%%%%%%%%%%%%%%%%%%%%%%%%%%%%%%%%%%%%%%%%%%%%%%%
\section{Annexes :}
TODO




\end{document}